\noindent\textbf{(a)} $V=50+5t$, $S=10t$ para $0\le t\le3$; $V=68-t$, $S=2(68-t)-\dfrac{100}{65^6}(68-t)^6$ para $3\le t<68$. V se mide en litros y S en gramos.

\noindent\textbf{(b)} $c=\dfrac{2t}{10+t}$ g/L para $0\le t\le3$ y $c=2-\dfrac{100}{65^6}(68-t)^5$ g/L para $3\le t<68$. Cuando $t\to68^-$, $S\to0$ y $c\to2$ g/L; el modelo no define concentración ni evolución posterior en un tanque vacío.

\noindent\textbf{(c)} $V_*=65(13/60)^{1/5}\approx47.8709$ L; $t_*=68-V_*\approx20.1291$ min; $S_*=5V_*/3\approx79.7849$ g.
